% !TeX root = ../analysis1.tex
\chapter{Differentiation}\label{chap:analysis-differentiation}

\section{Derivatives and linear approximation}
For the line $y=ax+b$, the slope is $(y_2-y_1)/(x_2-x_1)=a$.
For a general graph, consider the secant through $(c,f(c))$ and
$(c+h,f(c+h))$. Its slope is $[f(c+h)-f(c)]/h$.
The derivative is the limiting slope as the second point approaches the first.

\begin{definition}[Derivative]
Let $f$ be defined on an open interval containing $c$. It is differentiable
at $c$ if the finite limit
\[
f'(c)=\lim_{h\to0}\frac{f(c+h)-f(c)}h
=\lim_{x\to c}\frac{f(x)-f(c)}{x-c}
\]
exists. Differentiability on an open interval means differentiability at
every point there. Another notation is $\left.\frac{df}{dx}\right|_{x=c}$.
\end{definition}

\begin{example}[Derivatives from the definition]
For $f(x)=p$ constant, the difference quotient is zero, so $f'=0$.
For $f(x)=x$, it equals $1$, so $f'=1$.
As an additional example, for $f(x)=x^2$ it equals
\[
\frac{(c+h)^2-c^2}{h}=2c+h\longrightarrow2c.
\]
The quadratic calculation replaces the source's incomplete third example.
\end{example}

\begin{theorem}[Linear approximation]\label{thm:analysis-linear}
The function $f$ is differentiable at $c$ if and only if some $\mu\in\mathbb R$
and remainder $R(h)$ satisfy, for sufficiently small $h$,
\[
f(c+h)=f(c)+\mu h+hR(h),\qquad R(h)\longrightarrow0\quad(h\to0).
\]
In that case $\mu=f'(c)$.
\end{theorem}
\begin{proof}
If $f'(c)$ exists, put $\mu=f'(c)$ and
$R(h)=[f(c+h)-f(c)]/h-\mu$ for $h\ne0$, with $R(0)=0$.
Then the identity holds and $R(h)\to0$.
Conversely, subtract $f(c)$ and divide by $h\ne0$ in the asserted identity.
The quotient is $\mu+R(h)\to\mu$, proving differentiability.
\end{proof}
The affine expression $f(c)+f'(c)h$ is the linear approximation. Its error
$hR(h)$ is smaller than $h$ in the precise sense that the ratio tends to
zero. The condition $R(h)\to0$ is essential.

\begin{theorem}[Differentiability implies continuity]
If $f$ is differentiable at $c$, it is continuous at $c$.
\end{theorem}
\begin{proof}
Linear approximation gives $f(c+h)-f(c)=h(f'(c)+R(h))\to0$.
\end{proof}
The converse fails: $|x|$ is continuous at $0$, but its left and right
difference quotients are $-1$ and $1$.

\section{Rules for differentiation}
\begin{theorem}[Sum, constant multiple and product rules]
If $f,g$ are differentiable at $c$ and $a\in\mathbb R$, then
\begin{align*}
(f+g)'(c)&=f'(c)+g'(c),\\
(af)'(c)&=af'(c),\\
(fg)'(c)&=f'(c)g(c)+f(c)g'(c).
\end{align*}
\end{theorem}
\begin{proof}
For sums and constant multiples, expand the difference quotients and use
the algebra of limits. For products, add and subtract $f(c)g(c+h)$:
\begin{align*}
\frac{f(c+h)g(c+h)-f(c)g(c)}h
&=\frac{f(c+h)-f(c)}h\,g(c+h)\\
&\quad+f(c)\frac{g(c+h)-g(c)}h.
\end{align*}
Continuity gives $g(c+h)\to g(c)$, so passing to limits gives the product
rule. This shortens the source's equivalent remainder expansion.
\end{proof}

\begin{theorem}[Sandwich theorem for derivatives]
Suppose $f,q$ are differentiable at $c$, $f(c)=q(c)=L$, and
$f'(c)=q'(c)=m$. If $f(x)\leq g(x)\leq q(x)$ near $c$, including at $c$,
then $g$ is differentiable at $c$ and $g'(c)=m$.
\end{theorem}
\begin{proof}
First $g(c)=L$. For $h>0$, subtract $L$ and divide by $h$ to sandwich the
difference quotient for $g$ between those for $f$ and $q$. Both tend to $m$.
For $h<0$, the inequalities reverse, but both outer limits are still $m$.
Thus both one-sided limits of the middle quotient equal $m$.
\end{proof}

\begin{theorem}[Chain rule]
If $g$ is differentiable at $c$ and $f$ is differentiable at $g(c)$,
with the composition defined near $c$, then
\[
(f\circ g)'(c)=f'(g(c))g'(c).
\]
\end{theorem}
\begin{proof}
By linear approximation,
\[
g(c+h)-g(c)=h(g'(c)+R_g(h))=:k(h),\qquad R_g(h)\to0.
\]
Hence $k(h)\to0$. Also
\[
f(g(c)+k)=f(g(c))+f'(g(c))k+kR_f(k),\qquad R_f(k)\to0.
\]
Set $R_f(0)=0$ so substitution remains valid even when $k(h)=0$.
For $h\ne0$, the composite difference quotient becomes
\[
\frac{f(g(c+h))-f(g(c))}h
=(g'(c)+R_g(h))\bigl(f'(g(c))+R_f(k(h))\bigr).
\]
Passing to the limit proves the formula without dividing by $k(h)$.
\end{proof}

\section{Powers, exponentials and inverse functions}
\begin{example}[Integer powers]
Induction and the product rule give $(x^n)'=nx^{n-1}$ for $n\in\mathbb N$.
The base case is $(x)'=1$. If it holds for $n$, then
\[
(x^{n+1})'=(xx^n)'=x^n+x(nx^{n-1})=(n+1)x^n.
\]
For $x\ne0$, the reciprocal has difference quotient
\[
\frac{1/(x+h)-1/x}{h}=-\frac1{x(x+h)}\longrightarrow-\frac1{x^2}.
\]
The chain rule applied to $1/x^{-n}$ gives the same power rule for negative
integer $n$, on $x\ne0$. The derivative of $x^0=1$ is zero.
\end{example}

\begin{proposition}[Derivative of the exponential]
For every $x\in\mathbb R$, $(\exp x)'=\exp x$.
\end{proposition}
\begin{proof}
For $|h|<1$, Lemma~\ref{lem:analysis-exp-bounds} gives
$1+h\leq\exp(h)\leq1/(1-h)$. Both bounding functions equal $1$ at $0$
and have derivative $1$ there. The derivative sandwich theorem gives
$\exp'(0)=1$. At arbitrary $x$, the addition law yields
\[
\frac{\exp(x+h)-\exp(x)}h
=\exp(x)\frac{\exp(h)-1}h\longrightarrow\exp(x).
\]
\end{proof}
\begin{corollary}
For $a>0$, $\frac d{dx}a^x=a^x\log a$.
\end{corollary}
\begin{proof}
Use $a^x=\exp(x\log a)$ and the chain rule.
\end{proof}

\begin{proposition}[Derivative of an inverse: conditional formula]
If $f$ is bijective and differentiable at $c$, and $f^{-1}$ is differentiable
at $f(c)$, then $f'(c)\ne0$ and
\[
(f^{-1})'(f(c))=\frac1{f'(c)}.
\]
\end{proposition}
\begin{proof}
Differentiate $f^{-1}(f(x))=x$ at $c$. The chain rule gives
$(f^{-1})'(f(c))f'(c)=1$.
\end{proof}
This calculation assumes differentiability of the inverse. The inverse
function theorem at the end of the chapter supplies sufficient conditions.

\begin{example}[Logarithms and real powers]
Applying the inverse function theorem to $\exp$, whose derivative is
positive, justifies differentiability of $\log$. Consequently, for $x>0$,
\[
\frac d{dx}\log x=\frac1{\exp(\log x)}=\frac1x.
\]
For $\alpha\in\mathbb R$, $x^\alpha=\exp(\alpha\log x)$ on $x>0$, so
\[
\frac d{dx}x^\alpha=x^\alpha\frac\alpha x=\alpha x^{\alpha-1}.
\]
For arbitrary real $\alpha$, the restriction $x>0$ matters. For instance,
$\sqrt x$ has no finite right derivative at $0$ because
$\sqrt h/h=1/\sqrt h\to+\infty$ as $h\to0^+$.
\end{example}

\section{Extrema and stationary points}
\begin{definition}[Global and local extrema]
A global maximum of $f:S\to\mathbb R$ at $c\in S$ means that $f(x)\leq f(c)$
for all $x\in S$. A global minimum reverses the inequality.
A local maximum means that the inequality holds for all domain points
with $|x-c|<\delta$, for some $\delta>0$. A local minimum is analogous.
An extremum is a maximum or a minimum.
\end{definition}
For two-sided derivative tests, $c$ must be an interior point of the domain.
A stationary point is one where $f'(c)=0$, called a critical point in the
source. Some texts use ``critical point'' more broadly to include points
where the derivative is undefined.

\begin{theorem}[Fermat's stationary-point theorem]
If $f$ is differentiable at an interior point $c$ where it has a local
extremum, then $f'(c)=0$.
\end{theorem}
\begin{proof}
At a local maximum, $f(c+h)-f(c)\leq0$ for small $h$.
For $h>0$, division and the limit give $f'(c)\leq0$.
For $h<0$, they give $f'(c)\geq0$. Hence $f'(c)=0$.
Apply the result to $-f$ for a minimum.
\end{proof}
The converse fails: $x^3$ has derivative zero at $0$ but no extremum there.
Endpoint extrema need not have derivative zero: $f(x)=x$ on $[-1,1]$ has
its minimum at $-1$ and maximum at $1$. To find extrema on a closed interval,
compare interior candidates with endpoint values.

\section{Rolle's theorem and the mean value theorem}
\begin{theorem}[Rolle's theorem]
If $f:[a,b]\to\mathbb R$ is continuous on $[a,b]$, differentiable on
$(a,b)$, and $f(a)=f(b)$, with $a<b$, then some $c\in(a,b)$ has $f'(c)=0$.
\end{theorem}
\begin{proof}
By the extreme value theorem, $f$ attains a maximum and minimum.
If the values agree, $f$ is constant. Otherwise at least one extreme
value differs from the common endpoint value, so is attained at an
interior point. Its derivative is zero by Fermat's theorem.
\end{proof}

\begin{theorem}[Mean value theorem]\label{thm:analysis-mvt}
If $f:[a,b]\to\mathbb R$ is continuous on $[a,b]$ and differentiable on
$(a,b)$, some $c\in(a,b)$ satisfies
\[
f'(c)=\frac{f(b)-f(a)}{b-a}.
\]
\end{theorem}
\begin{proof}
Put $m=[f(b)-f(a)]/(b-a)$ and $g(x)=f(x)-mx$.
Then $g(a)=g(b)$, so Rolle's theorem gives $g'(c)=0$ for some $c\in(a,b)$.
Since $g'=f'-m$, the result follows.
\end{proof}

\begin{corollary}[Derivative sign and monotonicity]
Under the mean value theorem hypotheses, $f'>0$ throughout $(a,b)$ implies
$f$ is strictly increasing on $[a,b]$, and $f'<0$ implies it is strictly
decreasing. If $f'=0$ throughout $(a,b)$, $f$ is constant.
Non-strict derivative inequalities give non-strict monotonicity.
\end{corollary}
\begin{proof}
For $a\leq x<y\leq b$, apply the theorem on $[x,y]$ to get
$f(y)-f(x)=f'(c)(y-x)$. The positive factor $y-x$ preserves the relevant sign.
\end{proof}

\section{First, second and higher derivative tests}
\begin{theorem}[First derivative test]
Let $f$ be continuous near $c$ and differentiable on each side.
If for some $\delta>0$,
\[
f'(x)>0\ (c-\delta<x<c),\qquad f'(x)<0\ (c<x<c+\delta),
\]
then $c$ is a strict local maximum. Reversing the two signs gives a
strict local minimum.
\end{theorem}
\begin{proof}
The mean value theorem on $[x,c]$ for $x<c$ gives $f(c)>f(x)$;
on $[c,x]$ for $x>c$ it gives $f(x)<f(c)$. Thus $c$ is a strict maximum.
Apply this to $-f$ for the minimum case.
\end{proof}

\begin{theorem}[Second derivative test]
Let $f$ be twice continuously differentiable near $c$, with $f'(c)=0$.
If $f''(c)>0$, then $c$ is a strict local minimum. If $f''(c)<0$, then
$c$ is a strict local maximum. If $f''(c)=0$, the test is inconclusive.
\end{theorem}
\begin{proof}
If $f''(c)>0$, continuity keeps $f''$ positive near $c$. Thus $f'$ is
strictly increasing there and, since $f'(c)=0$, is negative on the left
and positive on the right. Use the first derivative test.
Apply this to $-f$ when $f''(c)<0$.
\end{proof}

\begin{theorem}[Higher derivative test]
Suppose $f$ is $n$ times continuously differentiable near $c$, $n\geq2$,
with $f'(c)=\cdots=f^{(n-1)}(c)=0$ and $f^{(n)}(c)\ne0$.
For odd $n$, $c$ is not a local extremum.
For even $n$, $c$ is a strict local minimum if $f^{(n)}(c)>0$ and a
strict local maximum if $f^{(n)}(c)<0$.
\end{theorem}
\begin{proof}[Supplementary proof]
Assume $f^{(n)}(c)>0$. It stays positive near $c$.
Repeated applications of the mean value theorem, using the vanishing
lower derivatives at $c$, show that to the right of $c$ all lower
positive-order derivatives are positive and $f(x)-f(c)>0$.
To the left, the signs alternate: $f^{(n-1)}(x)<0$,
$f^{(n-2)}(x)>0$, and ultimately $f(x)-f(c)$ has sign $(-1)^n$.
For even $n$, both sides are above $f(c)$; for odd $n$ they are on opposite
sides. The negative case follows by applying this reasoning to $-f$.
\end{proof}
The source marks this result non-examinable for its 2024--25 course and
omits the proof.

\begin{example}[Vanishing derivatives]
For $f(x)=x^3$, $f'(0)=f''(0)=0$ and $f'''(0)=6$: there is no extremum at
$0$. By contrast, $x^4$ has a strict minimum there, and $-x^4$ has a strict
maximum. Thus $f''(c)=0$ does not decide the issue.

The flat function
\[
F(x)=\begin{cases}e^{-1/x^2},&x\ne0,\\0,&x=0\end{cases}
\]
has $F^{(k)}(0)=0$ for every $k\geq0$, yet $F(x)>0$ for $x\ne0$.
It has a strict minimum although no first nonzero derivative exists.

\textit{Supplementary explanation.} For every integer $m\geq0$,
$|x|^{-m}e^{-1/x^2}\to0$: put $t=1/x^2$ and use $e^t\geq t^j/j!$ for
an integer $j>m/2$. Repeated differentiation away from zero produces
$P_k(1/x)e^{-1/x^2}$ for some polynomial $P_k$. These expressions and their
difference quotients at zero tend to zero by that estimate. Induction
proves smoothness and all stated derivatives. The one-sided flat function
from Chapter~\ref{chap:analysis-foundations} has the same property at zero.
\end{example}

\section{The inverse function theorem}
This theorem gives a sufficient condition for a differentiable local
inverse. A zero derivative does not preclude invertibility: $x^3$ is
bijective on $\mathbb R$, although its derivative at zero is zero.
Its inverse $x^{1/3}$ has no finite derivative there.

\begin{theorem}[One-dimensional inverse function theorem]
Let $f$ be continuously differentiable on an open interval about $x_0$,
with $f'(x_0)\ne0$. Some open interval $J$ about $x_0$ is mapped bijectively
onto an open interval $I$ about $f(x_0)$. The inverse is differentiable,
and for all $y\in I$,
\[
(f^{-1})'(y)=\frac1{f'(f^{-1}(y))}.
\]
In particular, $(f^{-1})'(f(x_0))=1/f'(x_0)$.
\end{theorem}
\begin{proof}[Supplementary proof]
Suppose $f'(x_0)>0$; otherwise apply the argument to $-f$.
Continuity gives a neighbourhood where $f'>0$. Choose $u<x_0<v$ so that
$[u,v]$ lies in that neighbourhood. The mean value theorem makes $f$
strictly increasing there. By the intermediate value theorem it maps
$J=(u,v)$ bijectively onto $I=(f(u),f(v))$.

Let $g=f^{-1}$. To prove its continuity at $y\in I$, set $x=g(y)$.
For every sufficiently small $\varepsilon>0$, strict increase gives
$f(x-\varepsilon)<y<f(x+\varepsilon)$. All $z\in I$ sufficiently close
to $y$ lie between those two values, so
$x-\varepsilon<g(z)<x+\varepsilon$. Hence $g(z)\to g(y)$.
For $z\ne y$, put $x'=g(z)\ne x$. Then $x'\to x$ and
\[
\frac{g(z)-g(y)}{z-y}
=\left(\frac{f(x')-f(x)}{x'-x}\right)^{-1}
\longrightarrow\frac1{f'(x)}.
\]
This proves the claimed differentiability and formula.
\end{proof}
The source marks this material non-examinable for its 2024--25 course
and leaves its proof incomplete. The added proof avoids assuming the
inverse differentiable before that fact is established.
